% Part: first-order-logic
% Chapter: models-theories
% Section: set-theory

\documentclass[../../../include/open-logic-section]{subfiles}

\begin{document}

\olfileid[pt]{fol}{mat}{set}

\olsection{A Teoria dos Conjuntos}

Quase toda a matemática pode ser desenvolvida na teoria dos conjuntos.
Desenvolver a matemática nessa teoria envolve uma série de coisas.
Primeiro, requer um conjunto de axiomas para a relação~$\in$.  Vários
sistemas de axiomas diferentes foram desenvolvidos, às vezes com
propriedades conflitantes de~$\in$.  O sistema de axiomas conhecido como
$\Log{ZFC}$, a teoria dos conjuntos de Zermelo--Fraenkel com o axioma da escolha,
destaca-se: é de longe o mais amplamente usado e estudado, porque
acontece que seus axiomas bastam para provar quase todas as coisas
que os matemáticos esperam ser capazes de provar.  Mas antes que isso possa ser
estabelecido, é primeiro necessário deixar claro como podemos sequer
\emph{expressar} todas as coisas que os matemáticos gostariam de expressar.
Para começar, a linguagem não contém !!{constant}s nem !!{function}s,
então parece, à primeira vista, pouco claro que possamos falar sobre conjuntos
particulares (como $\emptyset$ ou $\Nat$), falar sobre operações sobre
conjuntos (como $X \cup Y$ e $\Pow{X}$), muito menos outras
construções que envolvem coisas diferentes de conjuntos, como relações
e funções.

Para começar, ``é um elemento de'' não é a única relação na qual estamos
interessados: ``é um subconjunto de'' parece quase tão importante.  Mas nós
podemos \emph{definir} ``é um subconjunto de'' em termos de ``é um elemento de''.
Para fazer isso, temos que encontrar !!a{formula}~$!A(x, y)$ na
linguagem da teoria dos conjuntos que é satisfeita por um par de
conjuntos~$\tuple{X, Y}$ sse $X \subseteq Y$.  Mas $X$ é um subconjunto de $Y$
exatamente quando todos os !!{element}s de~$X$ são também !!{element}s de~$Y$.  Então
podemos definir $\subseteq$ pela fórmula
\[
\lforall[z][(z \in x \lif z \in y)]
\]
Agora, sempre que quisermos usar a relação~$\subseteq$ em uma fórmula, nós
poderíamos, em vez disso, usar aquela fórmula (com $x$ e $y$ adequadamente substituídos,
e a variável ligada~$z$ renomeada se necessário).  Por exemplo,
a extensionalidade dos conjuntos significa que se quaisquer conjuntos~$x$ e $y$ estão
contidos um no outro, então $x$ e $y$ devem ser o mesmo conjunto. Isso
pode ser expresso por $\lforall[x][\lforall[y][((x \subseteq y \land y
    \subseteq x) \lif x = y)]]$, ou, se substituirmos $\subseteq$ pela
definição acima, por
\[
\lforall[x][\lforall[y][((\lforall[z][(z \in x \lif z \in y)] \land
    \lforall[z][(z \in y \lif z \in x)]) \lif x = y)]].
\]
Este é, de fato, um dos axiomas de $\Log{ZFC}$, o ``axioma da
extensionalidade''.

Não há !!{constant} para $\emptyset$, mas podemos expressar ``$x$ é
vazio'' por $\lnot \lexists[y][y \in x]$.  Então ``$\emptyset$ existe''
torna-se a !!{sentence}~$\lexists[x][\lnot \lexists[y][y \in
    x]]$. Este é outro axioma de~$\Log{ZFC}$.  (Note que o axioma
da extensionalidade implica que há apenas um conjunto vazio.)  Sempre que
quisermos falar sobre $\emptyset$ na linguagem da teoria dos conjuntos, nós
escreveríamos isso como ``há um conjunto que é vazio e \dots'' Como
exemplo, para expressar o fato de que $\emptyset$ é um subconjunto de todo
conjunto, poderíamos escrever
\[
\lexists[x][(\lnot\lexists[y][y \in x] \land \lforall[z][x \subseteq
    z])]
\]
onde, é claro, $x \subseteq z$ teria, por sua vez, que ser substituído por
sua definição.

Para falar sobre operações sobre conjuntos, como $X \cup Y$ e $\Pow{X}$, nós
temos que usar um truque semelhante.  Não há símbolos de função na
linguagem da teoria dos conjuntos, mas podemos expressar as relações funcionais $X
\cup Y = Z$ e $\Pow{X} = Y$ por
\begin{align*}
& \lforall[u][((u \in x \lor u \in y) \liff u \in z)]\\
& \lforall[u][(u \subseteq x \liff u \in y )]
\end{align*}
já que os !!{element}s de $X \cup Y$ são exatamente os conjuntos que são
ou !!{element}s de~$X$ ou !!{element}s de~$Y$, e os
!!{element}s de $\Pow{X}$ são exatamente os subconjuntos de~$X$.  No entanto,
isso não nos permite usar $x \cup y$ ou $\Pow{x}$ como se fossem
termos: só podemos usar as !!{formula}s inteiras que definem as
relações $X \cup Y = Z$ e $\Pow{X} = Y$. De fato, não sabemos
que essas relações são alguma vez satisfeitas, isto é, não sabemos que
uniões e conjuntos potência sempre existem. Por exemplo, a !!{sentence}
$\lforall[x][\lexists[y][\Pow{x} = y]]$ é outro axioma
de~$\Log{ZFC}$ (o axioma do conjunto potência).

E quanto a falar de pares ordenados ou funções?  Aqui temos que
explicar como podemos pensar em pares ordenados e funções como tipos especiais
de conjuntos.  Uma maneira de definir o par ordenado $\tuple{x, y}$ é
como o conjunto $\{\{x\}, \{x, y\}\}$.  Mas, como antes, não podemos
introduzir !!a{function} que nomeie esse conjunto; só podemos definir a
relação $\tuple{x, y} = z$, isto é, $\{\{x\}, \{x, y\}\} = z$:
\[
\lforall[u][(u \in z \liff (\lforall[v][(v \in u \liff v = x)] \lor
  \lforall[v][(v \in u \liff (v = x \lor v = y))]))]
\]
Isso diz que os !!{element}s~$u$ de~$z$ são exatamente aqueles conjuntos que
ou têm $x$ como seu único !!{element} ou têm $x$ e~$y$ como seus
únicos !!{element}s (em outras palavras, aqueles conjuntos que são ou idênticos
a $\{x\}$ ou idênticos a $\{x, y\}$).  Uma vez que temos isso, podemos dizer
outras coisas, por exemplo, que $X \times Y = Z$:
\[
\lforall[z][(z \in Z \liff \lexists[x][\lexists[y][(x \in
        X \land y \in Y \land \tuple{x, y} = z)]])]
\]

Uma função $f \colon X \to Y$ pode ser pensada como a relação $f(x)
= y$, isto é, como o conjunto de pares~$\Setabs{\tuple{x,y}}{f(x) = y}$. Nós
podemos então dizer que um conjunto~$f$ é uma função de $X$ em $Y$ se (a) é
uma relação $\subseteq X \times Y$, (b) é total, isto é, para todo $x
\in X$ há algum $y \in Y$ tal que $\tuple{x, y} \in f$ e (c)
é funcional, isto é, sempre que $\tuple{x, y}, \tuple{x, y'} \in f$,
$y = y'$ (porque os valores de funções devem ser únicos). Então ``$f$ é uma
função de $X$ em $Y$'' pode ser escrito como:
\begin{align*}
 \lforall[u][(u \in f \lif {}] & \lexists[x][\lexists[y][(x \in X \land y \in
      Y \land \tuple{x, y} = u)]]) \land {}\\
 \lforall[x][(x \in X \lif {}] &
      (\lexists[y][(y \in Y \land \mathrm{mapeia}(f, x, y))] \land {}\\
& (\lforall[y][\lforall[y'][((\mathrm{mapeia}(f, x, y) \land
    \mathrm{mapeia}(f, x, y')) \lif y = y')]]))
\end{align*}
onde $\mathrm{mapeia}(f, x, y)$ abrevia $\lexists[v][(v \in f \land
  \tuple{x, y} = v)]$ (esta !!{formula} expressa ``$f(x) = y$'').

Agora também não é difícil expressar que $f\colon X \to Y$ é
!!{injective}, por exemplo:
\begin{multline*}
 f \colon X \to Y \land \lforall[x][\lforall[x'][((x \in X \land x' \in
     X \land {}]] \\
 \lexists[y][(\mathrm{mapeia}(f, x, y) \land \mathrm{mapeia}(f,
        x', y))]) \lif x = x')
\end{multline*}
Uma função~$f\colon X \to Y$ é !!{injective} sse, sempre que $f$ mapeia $x, x'
\in X$ para um único~$y$, $x = x'$.  Se abreviarmos esta fórmula como
$\mathrm{inj}(f, X, Y)$, já estamos em posição de afirmar na
linguagem da teoria dos conjuntos algo tão não trivial quanto o teorema de Cantor:
não há função !!{injective} de $\Pow{X}$ em $X$:
\[
\lforall[X][\lforall[Y][(\Pow{X} = Y \lif
    \lnot\lexists[f][\mathrm{inj}(f, Y, X)])]]
\]

Poderíamos pensar que a teoria dos conjuntos requer outro axioma que garanta
a existência de um conjunto para toda propriedade definidora. Se $!A(x)$ é uma
fórmula da teoria dos conjuntos com a variável~$x$ livre, podemos considerar a
!!{sentence}
\[
\lexists[y][\lforall[x][(x \in y \liff !A(x))]].
\]
Esta !!{sentence} afirma que há um conjunto~$y$ cujos !!{element}s
são todos e somente aqueles $x$ que satisfazem~$!A(x)$. Este esquema é chamado de
``princípio da compreensão''.  Parece muito útil; infelizmente
é inconsistente.  Tome $!A(x) \ident \lnot x \in x$, então o
princípio da compreensão afirma
\[
\lexists[y][\lforall[x][(x \in y \liff x \notin x)]],
\]
isto é, afirma a existência de um conjunto de todos os conjuntos que não são
!!{element}s de si mesmos. Nenhum tal conjunto pode existir---este é o
Paradoxo de Russell.  $\Log{ZFC}$, de fato, contém uma versão restrita---e
consistente---deste princípio, o princípio da separação:
\[
\lforall[z][\lexists[y][\lforall[x][(x \in y \liff (x \in z \land
      !A(x))]]].
\]

\begin{prob}
Mostre que o princípio da compreensão é inconsistente dando
!!a{derivation} que mostre
\[
\lexists[y][\lforall[x][(x \in y \liff x \notin x)]] \Proves \lfalse.
\]
Pode ajudar mostrar primeiro $(A \lif \lnot A) \land (\lnot A \lif A)
\Proves \lfalse$.
\end{prob}
\end{document}
